
\section{Introduction}

In 2024, German day-ahead electricity prices turned negative for 459 hours, triple the figure recorded just two years earlier \cite{FfE2025_GermanElectricityPricesEPEX2024}. This is not a market anomaly. It is the structural consequence of deploying wind energy at scale: the more of it a power system absorbs, the less each additional unit is worth on the wholesale market. This tension between deployment ambition and economic erosion is the central problem this thesis investigates.

\vspace{4pt}

By the end of 2024, Europe had roughly 35 to 37~GW of installed offshore wind capacity \cite{Colaluce2025_EuropeOffshoreWind}, and long-term decarbonisation pathways call for approximately 360~GW by 2050 \cite{EuropeanCommission2021_EUReferenceScenario2020}. Most of this capacity sits in the North Sea, the world's most mature offshore basin, where Germany alone operates 9.2~GW \cite{OffshoreWIND2025_GermanyOffshoreCapacity}. The Baltic Sea remains smaller but is accelerating, with the Marienborg Declaration targeting 19.6~GW by 2030 and an estimated technical potential of up to 93~GW \cite{BalticWind2024_BEMIPOffshorePriorities}. The question is therefore no longer whether offshore wind will expand, but where and under what market conditions.

\vspace{4pt}

The political ambition behind this expansion is substantial and accelerating. The REPowerEU plan raised the EU's 2030 renewable energy target to 42.5\% of final energy consumption \cite{REPowerEU2022}, and both the IEA and IRENA project that global renewable capacity must at least triple by 2030 to meet stated climate goals \cite{IEA2024, IRENA2024}. Europe's competitive position adds a further imperative: the Draghi report on European competitiveness found that EU electricity prices are two to three times higher than those in the United States, identifying energy cost reform as a prerequisite for industrial survival \cite{Draghi2024}. The 2024 EU Electricity Market Design reform \cite{EUMarketReform2024} responded to these pressures by mandating two-sided contracts for difference for all new publicly supported renewable capacity, signalling that revenue stabilisation is now a structural feature of European energy policy rather than a transitional support measure. For offshore wind, these dynamics mean that the volume of zero-marginal-cost generation entering European markets will grow by an order of magnitude over the coming decades, amplifying the value erosion problem this thesis examines.

\vspace{4pt}

Under these conditions, project viability depends not on energy yield or levelised cost alone but on the price a wind farm actually receives in the hours it generates \cite{Hirth2013}, a metric that declines as penetration grows.

\vspace{4pt}

Substantial bodies of literature address wind resource assessment, electricity market dynamics, and spatial modelling, yet these remain largely disconnected. Wind resource studies quantify atmospheric regimes, capacity factors and spatial correlation, yet rarely link these characteristics to market value. Electricity market studies analyse merit order effects and renewable penetration, but typically at the national level and without detailed treatment of basin-level wind heterogeneity. GIS-based wind potential assessments classify sites by wind speed, depth and distance to shore, but largely omit capture value. No existing study systematically compares the North Sea and Baltic Sea, two basins with distinct wind regimes, generation mixes and interconnection structures, on the dimension that ultimately determines project viability: revenue.

\vspace{4pt}

This literature review synthesises the state of knowledge across these largely disconnected domains and critically evaluates existing modelling approaches in order to identify the gap this thesis addresses. The central research question emerges from the intersection of resource assessment, market economics, and spatial modelling: can geographical diversification of wind output improve merchant revenues (the income earned by selling directly on the wholesale market, without a guaranteed support contract) in the Baltic Sea relative to the North Sea, and can a transferable modelling framework be developed to compare offshore basins systematically?

\vspace{4pt}

The review proceeds through the economic mechanisms driving value erosion (Section~\ref{sec:value_problem}), the regional factors that cause this erosion to differ between the two basins (Section~\ref{sec:regional_drivers}), and the modelling tools available to quantify wind energy potential and market value across space (Section~\ref{sec:modelling}), before identifying the gaps in the existing literature and setting out the objectives of this thesis (Section~\ref{sec:research_gap}).
